% v12 Design: cascade, labels, safety layer, issue states.
\section{Design}\label{sec:design}

% Extracted from v11 paper-v11.tex by tools/v12/make_v12.py; do not edit.
\begin{figure*}[t]
\centering
\begin{tikzpicture}[
  node distance=4mm and 5mm, >={Stealth[length=2.2mm]}, font=\small,
  box/.style={draw=ink!70, rounded corners=2pt, fill=ink!6, align=center, inner sep=3pt, minimum height=9mm},
  tier/.style={draw=ink!80, rounded corners=2pt, fill=white, align=left, inner sep=3pt, text width=33mm, minimum height=9mm},
  lbl/.style={rounded corners=2pt, align=center, inner sep=2.5pt, minimum width=17mm, font=\small\sffamily},
  st/.style={draw=black!45, rounded corners=2pt, fill=black!4, align=center, inner sep=3pt, text width=27mm}]
\node[box, text width=29mm] (old) {Issue on version $t$\\ anchors $a=(\pi,s,e,c,q)$};
\node[box, text width=29mm, below=6mm of old] (new) {Files of version $t{+}1$\\ previous text of $\pi$};
\node[tier, right=8mm of old, yshift=7mm] (t1) {1 same path, same hash};
\node[tier, below=2mm of t1] (t2) {2 same stripped lines in $\pi$, nearest to $s$};
\node[tier, below=2mm of t2] (t3) {3 best similar window near the diff-predicted line, if ${\ge}\,\theta$};
\node[tier, below=2mm of t3] (t4) {4 same stripped lines in another file};
\node[lbl, fill=ink, text=white, right=5mm of t1] (l1) {identical};
\node[lbl, fill=ink, text=white, right=5mm of t2, yshift=2.5mm] (l2a) {content-match};
\node[lbl, fill=cyan!25, below=1mm of l2a] (l2b) {moved};
\node[lbl, fill=orange!35, right=5mm of t3] (l3) {changed};
\node[lbl, fill=cyan!25, right=5mm of t4] (l4) {moved};
\node[lbl, fill=magenta!20, below=3mm of l4] (l5) {removed};
\node[st, right=13mm of l2a, yshift=-1mm] (keep) {\textbf{keeps its state}\\ \footnotesize prior active state; carried with lineage};
\node[st, right=13mm of l3, yshift=-3mm] (re) {\textbf{needs re-check}\\ \footnotesize supervisor sees old quote next to new place};
\node[st, below=4mm of re, fill=white, draw=black!30, dashed] (res) {\textbf{resolved}\\ \footnotesize only by the supervisor};
\begin{scope}[on background layer]
\node[draw=ink!30, fill=ink!2, rounded corners=3pt, fit=(t1)(t4), inner sep=3pt, label={[font=\footnotesize\itshape, text=ink]above:relocation cascade (first match wins)}] {};
\end{scope}
\draw[->, ink] (old.east) -- ++(4mm,0) |- (t1.west);
\draw[->, ink] (new.east) -- ++(4mm,0) |- (t1.west);
\foreach \a/\b in {t1/t2, t2/t3, t3/t4} \draw[->, black!50] (\a.south) -- node[right, font=\scriptsize, black!60]{no} (\b.north);
\draw[->, black!50] (t4.south) |- node[pos=0.25, right, font=\scriptsize, black!60]{no} (l5.west);
\draw[->] (t1) -- (l1); \draw[->] (t2.east) -- (l2a.west); \draw[->] (t2.east) -- (l2b.west);
\draw[->] (t3) -- (l3); \draw[->] (t4) -- (l4);
\foreach \n in {l1,l2a,l2b,l4} \draw[->, black!55] (\n.east) -- (keep.west);
\foreach \n in {l3,l5} \draw[->, warn] (\n.east) -- (re.west);
\draw[->, black!40, dashed] (re.south) -- (res.north);
\end{tikzpicture}
\caption{Carry-forward of one issue. Each anchor is re-found by the first tier that matches; the issue needs a re-check if any anchor is \lbl{changed} or \lbl{removed}. No automated transition leads to \lbl{resolved}: only the supervisor closes an issue.}
\label{fig:cascade}
\end{figure*}


\subsection{Versions, Anchors and the Cascade}
A submission version $V_t$ contains files $(\pi,T)$ and extracted report text. A code anchor is $a=(\pi,s,e,c,q)$: path, first and last line of a span of height $k=e-s+1$, whole-file SHA-256 digest and quoted text. Let $\qbar$ be the nonblank quoted lines with surrounding whitespace removed. The quotation is matching evidence, not a unique identifier of a program entity. An issue may hold up to twelve anchors.

Relocation returns a new position and a label (Fig.~\ref{fig:cascade}). Tier~1 accepts an untouched file. Tier~2 searches $\qbar$ as a contiguous run of stripped non-blank lines in the same file and takes the occurrence nearest the old start. Tier~3 handles edits in place: it maps the old start through a line diff \cite{pythonsequence} of the previous and new file to an expected start $p$, scores every window of height $k$ starting within $p \pm r$, $r=\max(10,2k)$, by the similarity ratio $2\sum_i m_i/(|u|+|v|)$ of the joined stripped lines, where $m_i$ are the matching-block lengths, and accepts the best window if the ratio reaches $\theta$ (deployed default $0.85$). Tier~4 searches $\qbar$ in the other files in sorted path order. Tier~3 runs before tier~4 so that an edited block is recognized in place rather than matched to an identical copy elsewhere, and its window stays near the predicted line to avoid distant look-alikes.
% Extracted from v11 paper-v11.tex by tools/v12/make_v12.py; do not edit.
\begin{table}[t]
\caption{Relocation Labels: Definition, Limits and Effect on the Issue}\label{tab:labels}
\centering\footnotesize
\setlength{\tabcolsep}{3pt}
\begin{tabular}{@{}p{.17\columnwidth}p{.37\columnwidth}p{.24\columnwidth}p{.12\columnwidth}@{}}\toprule
Label & Holds when & Does \emph{not} establish & Issue \\ \midrule
\lbl{identical} & same path and same file digest (tier 1) & semantic review of the file & keeps state \\
\lbl{content-match} & normalized quotation $\qbar$ found at the old start (tier 2) & byte identity, equal indentation, equal behavior & keeps state \\
\lbl{moved} & $\qbar$ found at another start or in another file (tiers 2, 4) & that the match is the original copy & keeps state \\
\lbl{changed} & best near-predicted window has ratio $\ge\theta$ (tier 3) & that the concern is affected & re-check \\
\lbl{removed} & no tier succeeds & that the code was deleted & re-check \\ \bottomrule
\end{tabular}
\tablenote{}{The deployed code emits a single identifier, \texttt{exact}, for both \lbl{identical} and \lbl{content-match}; this paper separates them because only the first is a byte-level property. Normalization strips each line's leading and trailing whitespace and skips blank lines.}
\end{table}


\subsection{Labels and Their Limits}
Table~\ref{tab:labels} defines the five labels; \lbl{changed} and \lbl{removed} return the issue for re-check. Normalization lets an anchor survive blank-line and formatting edits, but in Python it also hides indentation, which carries meaning: a statement dedented out of a conditional keeps its stripped text while its behavior changes. A \lbl{content-match} or \lbl{moved} label is therefore textual, not semantic, evidence. A second limit concerns copies. When the original is edited or deleted and an identical copy exists elsewhere, tiers 2 and 4 can return the copy as \lbl{moved}, which keeps the issue's state, so the change is not returned to the supervisor.

\subsection{The Safety Layer}\label{sec:guards}
The layer consists of the labels, two window guards, an identity rule and the issue states. The guards were derived from the failure analysis of Section~\ref{sec:layer} and evaluated under protocols whose text and code hashes were recorded before any held-out run; the deployed cascade remains the system of record.
\begin{itemize}
\item \textbf{Indentation guard (IG).} When a normalized run is accepted, the leading-whitespace widths of its non-blank lines are compared with the quotation's, relative to the first line and in absolute base indentation. For Python any difference turns the label into \lbl{changed}; for other languages only a relative difference does.
\item \textbf{Copy guard (AGi).} When a run is accepted anywhere other than the predicted start, the guard counts occurrences of $\qbar$ in the old and new versions of the candidate file. A genuine move introduces the text, so only an increased count yields a clean \lbl{moved}; otherwise the placement is kept but re-check is requested. The first version of this guard (AG) flagged every clean insertion directly above an anchor, because the pinned diff maps such an anchor's old start to the first inserted line; AGi also accepts the line to which the diff maps the anchor's own first line as the predicted start.
\item \textbf{Variable-height window (VW).} Tier~3 also scores windows of heights $k\pm d$, $d=\max(1,\lceil k/2\rceil)$, so that statements that grow or shrink are recovered with their new extent. VW belongs to the cascade's own search.
\item \textbf{Identity rule (ID).} When tier~3 accepts a \lbl{changed} window, carry-forward replaces the anchor's quotation with the matched text, so a wrong fuzzy match becomes the reference for every later version. ID keeps the quotation the supervisor reviewed, $q_0$, beside the current one, $q_t$. Relocation still tracks $q_t$, but every \lbl{changed} hop must reach similarity 0.5 to $q_0$, otherwise the anchor is reported removed, and the request is computed against $q_0$: an issue returns whenever the located code differs from what the supervisor reviewed.
\end{itemize}
IG and AGi need only the returned window and the two versions of the file, and ID only the quotations, so all three apply to the output of any locator. None proves identity: two copies introduced in the same revision defeat AGi, and IG is a layout signal, not an equivalence test. A provenance rule (AG2), which accepts a window by where its lines came from in the diff, is a negative result (Section~\ref{sec:layer}). Content alone cannot decide which of two identical copies is the reviewed one; it can only decide when to ask.

% Extracted from v11 paper-v11.tex by tools/v12/make_v12.py; do not edit.
\begin{figure}[t]
\centering
% Review-issue lifecycle across versions (shared by both papers).
\begin{tikzpicture}[>={Stealth[length=2mm]}, font=\scriptsize, node distance=7mm and 4.5mm,
  st/.style={draw=ink!70, rounded corners=6pt, fill=ink!6, align=center, inner sep=3pt, minimum width=14mm, minimum height=6mm, font=\scriptsize\sffamily},
  fin/.style={st, fill=black!6, draw=black!50, double}]
\node[st] (open) {open};
\node[st, right=of open] (acc) {accepted};
\node[st, below=9mm of $(open)!0.5!(acc)$, fill=orange!30, draw=warn] (re) {needs re-check};
\node[fin, right=of acc] (res) {resolved};
\node[fin, below=of res] (dis) {dismissed};
\node[left=6mm of open, align=right, text=black!60] (new) {new finding\\ or faculty point};
\draw[->] (new) -- (open);
\draw[->] (open) -- (acc);
\draw[->] (acc) -- (res);
\draw[->] (re.east) -- (dis.west);
\draw[->, warn] (open.south) -- node[pos=0.6, left=6pt, text=warn, align=right]{carried; an anchor\\ \lab{changed}/\lab{removed}} (re);
\draw[->, warn] (acc.south) -- (re);
\draw[->] (re.north east) -- (res.south west);
\draw[->, black!55] (open) edge[loop above] node[above, align=center]{carried; all anchors\ \lab{identical},\\ \lab{content-match} or \lab{moved}} (open);
\draw[->, black!55, dashed] (dis.south) -- ++(0,-5mm) node[below, align=center, text=black!60]{same finding on version $t{+}1$\\ arrives dismissed};
\end{tikzpicture}
\caption{Lifecycle of a review issue across versions. Black arrows are faculty actions, orange arrows are carry-forward. Carry-forward moves an issue to \lab{needs re-check} when any anchor is \lab{changed} or \lab{removed}, and otherwise keeps its state. Supervisors may set any state from any state; no system or assistant action resolves an issue. A new finding matching an issue dismissed on the preceding version is created dismissed, and a deleted issue is not carried (Section~\ref{sec:statecheck}).}
\label{fig:states}
\end{figure}


\subsection{Issue States}
When a new version arrives, active issues (\lbl{open}, \lbl{accepted} or \lbl{needs re-check}) are copied into its review with relocated anchors and a lineage reference. Any \lbl{changed} or \lbl{removed} anchor requests re-check; otherwise the previous state is kept. No relocation transition resolves an issue. A finding raised again by analysis is matched to a carried issue by category, title and overlapping evidence, so a carried concern is not duplicated. A new finding is also compared with the findings dismissed on the preceding version by category and the first 200 normalized characters of its description; a match is created already dismissed. That key is a heuristic: findings sharing an opening paragraph are conflated and rewording evades it (Section~\ref{sec:statecheck}). Notes attached to carried issues, including voice, video and image notes, follow their issue; notes pinned to code are re-found by the cascade. Fig.~\ref{fig:states} summarizes the lifecycle.

\subsection{Three Properties and Five Safety Requirements}\label{sec:props}
Let anchor $a$ of issue $i$ be reviewed with quotation $q_0$, and let $\rho_t(a)$ be its true location in version $V_t$, or $\bot$ if the reviewed code is gone. Three properties are kept apart. \emph{Change}: $\delta_t(a)$ holds if $\rho_t(a)=\bot$ or the code at $\rho_t(a)$ differs from $q_0$; edits around the span are outside $\delta$ (Section~\ref{sec:suggest}). \emph{Anchor validity}: the returned window $w_t(a)$ is valid if $w_t(a)=\rho_t(a)$. \emph{Concern}: $\kappa_t(i)$ holds while the supervisor's concern is unresolved. The first two concern code and can be scored against a reference; $\kappa$ is known only to the supervisor, and a change is evidence that the concern may be affected, not that it is. The layer computes a request $r_t(a)$ and the lifecycle a state $\sigma_t(i)$, and should satisfy:
\begin{itemize}
\item[\textbf{S1}] no unsignaled change: $\delta_t(a)\Rightarrow r_t(a)$;
\item[\textbf{S2}] no silent misdirection: $w_t(a)\notin\{\bot,\rho_t(a)\}\Rightarrow r_t(a)$;
\item[\textbf{S3}] visible absence: $w_t(a)=\bot\Rightarrow r_t(a)$;
\item[\textbf{S4}] human closing: every transition into \lbl{resolved} or \lbl{dismissed} is a supervisor action;
\item[\textbf{S5}] persistence: an active issue that no supervisor closes is active in $V_{t+1}$, and requests a check there if any of its anchors satisfies S1's premise.
\end{itemize}
Each carried anchor is thus \emph{verified} (placed, no request), \emph{ambiguous} (placed with a request; the copy guard has found the identical copies and could offer them as candidates) or \emph{missing} (not placed, with a request). Only a verified anchor is shown without a request. Where S1 and S2 hold, that is under A1--A3 and for changes inside the quoted span, a verified anchor is the reviewed code, unchanged; ``verified'' is this conditional implication, not a runtime certificate, and Table~\ref{tab:counter} lists where the assumptions fail. S3 holds by construction, because \lbl{removed} requests a check. S1 holds when the request is computed against $q_0$ and (A1) the window is valid or absent and (A2) normalization preserves meaning; S2 needs (A3) that a wrong window differs from $q_0$. S4 is a property of code paths and S5 of the lifecycle; both are checked in Section~\ref{sec:statecheck}. Table~\ref{tab:counter} lists every failure mechanism found, the assumption it breaks, a minimal counterexample and what remains after the layer.

\begin{table*}[t]
\caption{Failure Mechanisms, the Requirement Each Violates, and What Remains}\label{tab:counter}
\centering\footnotesize
\setlength{\tabcolsep}{3pt}
\begin{tabular}{@{}p{.13\textwidth}p{.07\textwidth}p{.37\textwidth}p{.14\textwidth}p{.24\textwidth}@{}}\toprule
Mechanism & Breaks & Minimal counterexample ($V_0\rightarrow V_t$) & Addressed by & Residual \\ \midrule
Identical copy & S1, S2 (A3) & \texttt{a.py} and \texttt{b.py} hold the same retry loop; the reviewed loop in \texttt{a.py} is edited, and the unchanged copy in \texttt{b.py} is returned as \lbl{moved} & copy guard (AG, AGi) & copies adjacent to the original or added by the same revision (Section~\ref{sec:copysuite}); the guard asks but does not find the reviewed copy \\
Hidden indentation & S1 (A2) & \texttt{if ok:} / \texttt{\ \ send()} becomes \texttt{send()} at the outer level; stripped text is equal, so \lbl{content-match} & IG & a layout signal, not equivalence; relative only outside Python \\
Wrong fuzzy hop & S1 (A1) over versions & a look-alike window is accepted as \lbl{changed} at hop 3, carry replaces the quotation, and later hops compare against the wrong text & identity rule & choice of floor; some placements become lost \\
Locator verdict & S1 & a language model reports a region unchanged although its quoted lines were edited & labels from $q_0$ & none for S1; placement stays the locator's \\
Insertion around the quotation & none: outside $\delta$ & a suggestion wraps the reviewed call in a new guard clause; the quoted lines stand verbatim, so no label changes (Section~\ref{sec:suggest}) & partly AG & a check of the lines around the span \\
Deleted target placed on look-alike & validity only & the reviewed block is deleted and a similar block nearby passes $\theta$; the request fires but points at other code & none & a misdirected request \\
Issue deletion & S5 & a supervisor or the assistant deletes an active issue; it is not carried & stronger contract (Sec.~\ref{sec:statecheck}) & implemented: journaled undo only \\
Dismissal key collision & S4, S5 & a different finding with the same first 200 characters arrives \lbl{dismissed} & stronger contract & implemented: heuristic key \\
Uncomposed version & S5 & $V_2$ is never composed, so $V_3$ receives none of $V_1$'s issues & stronger contract & implemented: needs recomposition \\ \bottomrule
\end{tabular}
\end{table*}
